%%%%%%%%%%%%%%%%%%%%%%%%%%%%%%%%%%%%%%%%%
% Beamer Presentation
% LaTeX Template
% Version 1.0 (10/11/12)
%
% This template has been downloaded from:
% http://www.LaTeXTemplates.com
%
% License:
% CC BY-NC-SA 3.0 (http://creativecommons.org/licenses/by-nc-sa/3.0/)
%
%%%%%%%%%%%%%%%%%%%%%%%%%%%%%%%%%%%%%%%%%

%------------------------------------------------------------------------------------------------
% PACKAGES AND THEMES
%------------------------------------------------------------------------------------------------

\documentclass[table, xcolor = {dvipsnames}, 9pt]{beamer}
\usepackage{tikz}
\usetikzlibrary{calc}
\usetikzlibrary{positioning}
\usetikzlibrary{arrows.meta}
\usetikzlibrary{external}
\mode<presentation> {

% The Beamer class comes with a number of default slide themes
% which change the colors and layouts of slides. Below this is a list
% of all the themes, uncomment each in turn to see what they look like.

\usetheme{default}
%\usetheme{AnnArbor}
%\usetheme{Antibes}
%\usetheme{Bergen}
%\usetheme{Berkeley}
%\usetheme{Berlin}
%\usetheme{Boadilla}
%\usetheme{CambridgeUS}
%\usetheme{Copenhagen}
%\usetheme{Darmstadt}
%\usetheme{Dresden}
%\usetheme{Frankfurt}
%\usetheme{Goettingen}
%\usetheme{Hannover}
%\usetheme{Ilmenau}
%\usetheme{JuanLesPins}
%\usetheme{Luebeck}
%\usetheme{Madrid}
\usetheme{metropolis}
%\usetheme{Malmoe}
%\usetheme{Marburg}
%\usetheme{Montpellier}
%\usetheme{PaloAlto}
%\usetheme{Pittsburgh}
%\usetheme{Rochester}
%\usetheme{Singapore}
%\usetheme{Szeged}
%\usetheme{Warsaw}

% As well as themes, the Beamer class has a number of color themes
% for any slide theme. Uncomment each of these in turn to see how it
% changes the colors of your current slide theme.

%\usecolortheme{albatross}
%\usecolortheme{beaver}
%\usecolortheme{beetle}
%\usecolortheme{crane}
%\usecolortheme{dolphin}
%\usecolortheme{dove}
%\usecolortheme{fly}
%\usecolortheme{lily}
%\usecolortheme{orchid}
%\usecolortheme{rose}
\usecolortheme{seagull}
%\usecolortheme{seahorse}
%\usecolortheme{whale}
%\usecolortheme{wolverine}
\usefonttheme{professionalfonts}
%\setbeamertemplate{footline} % To remove the footer line in all slides uncomment this line
%\setbeamertemplate{footline}[page number] % To replace the footer line in all slides with a simple slide count uncomment this line

%\setbeamertemplate{navigation symbols}{} % To remove the navigation symbols from the bottom of all slides uncomment this line
}

\usepackage{graphicx} % Allows including images
\usepackage{booktabs} % Allows the use of \toprule, \midrule and \bottomrule in tables
\usepackage{float}
\usepackage{tikz}
\usepackage{multirow}
\usepackage{natbib}
\usepackage{hyperref}
\usepackage{diagbox}
\usepackage{makecell}
\usepackage{xparse}
\usepackage{subfig}
\usepackage{amsmath}
\usepackage{amsfonts,amsthm,amsmath,amssymb}
\usepackage{bbm}
\usepackage{bm}
\usepackage{empheq}
\usepackage{pgfplots}
\usepackage{animate}
\usepgfplotslibrary{colorbrewer}
\pgfplotsset{compat=1.17}

\newcommand\mybox[2][]{\tikz[overlay]\node[fill=lightgray,inner sep=2pt, anchor=text, rectangle, rounded corners=1mm,#1] {#2};\phantom{#2}}
\hypersetup{unicode=true,
            bookmarksnumbered=true,
            bookmarksopen=true,
            bookmarksopenlevel=2,
            breaklinks=false,
            pdfborder={0 0 1},
            hypertexnames=false,
            pdfstartview={XYZ null null 1}}
\usepackage{xcolor}
% R code listing style (matches the course's R-code slides)
\usepackage{listings}
\definecolor{codegreen}{rgb}{0,0.6,0}
\definecolor{codegray}{rgb}{0.5,0.5,0.5}
\definecolor{codepurple}{rgb}{0.58,0,0.82}
\definecolor{backcolour}{rgb}{0.95,0.95,0.92}
\lstdefinestyle{Rstyle}{
    language=R,
    backgroundcolor=\color{backcolour},
    commentstyle=\color{codegreen},
    keywordstyle=\color{magenta},
    stringstyle=\color{codepurple},
    basicstyle=\ttfamily\footnotesize,
    breakatwhitespace=false,
    breaklines=true,
    keepspaces=true,
    showspaces=false,
    showstringspaces=false,
    showtabs=false,
    tabsize=2
}
\newcommand\myheading[1]{%
  \par\bigskip
  {\Large\bfseries#1}\par\smallskip}
\newcommand\given[1][]{\:#1\vert\:}
\theoremstyle{plain}
\newtheorem{thm}{Theorem}
\newtheorem{prop}{Proposition\thisthmnumber}
\newtheorem{lem}{Lemma\thisthmnumber}
\newtheorem{cor}{Corollary}
\newtheorem{defin}{Definition}
\newtheorem{algo}{Algorithm}
\newcommand*\diff{\mathop{}\!\mathrm{d}}
\newcommand*\Diff[1]{\mathop{}\!\mathrm{d^#1}}
\newcommand{\bh}[1]{{\color{blue}{#1}}}
\newcommand{\mh}[1]{{\color{magenta}{#1}}}
\newcommand{\thisthmnumber}{}
\newcommand{\tikzmark}[1]{\tikz[baseline,remember picture] \coordinate (#1) {};}
\newcommand*{\QEDA}{\hfill\ensuremath{\blacksquare}}%
\newcommand*{\QEDB}{\hfill\ensuremath{\square}}%
\DeclareMathOperator{\E}{\rm{E}}
\DeclareMathOperator{\R}{\mathbb{R}}
\DeclareMathOperator{\N}{\mathbb{N}}
\DeclareMathOperator{\Var}{\rm{Var}}
\DeclareMathOperator{\Cov}{\rm{Cov}}
\DeclareMathOperator{\Supp}{\rm{Supp}}
\DeclareMathOperator{\e}{\rm{e}}
\DeclareMathOperator{\F}{\mathcal{F}}
\DeclareMathOperator{\Z}{\mathcal{Z}}
\DeclareMathOperator{\logit}{\rm{logit}}
\DeclareMathOperator{\indep}{{\perp\!\!\!\perp}}
\DeclareMathOperator{\rank}{rank}
\DeclareMathOperator*{\argmin}{arg\,min}
\DeclareMathOperator*{\argmax}{arg\,max}
%\DeclareMathOperator{\Pr}{\rm{Pr}}
%------------------------------------------------------------------------
% TITLE PAGE
%-----------------------------------------------------------------------
\pagestyle{empty}
\title[]{Uncertainty Estimation, Testing, and Confidence Intervals} % The short title appears at the bottom of every slide, the full title is only on the title page

\author{Thomas Leavitt} % Your name
\institute[] % Your institution as it will appear on the bottom of every slide, may be shorthand to save space
{
% Your institution for the title page
\medskip
\textit{} % Your email address
}
\date{June 18, 2026} % Date, can be changed to a custom date

\begin{document}

\begin{frame}
\titlepage % Print the title page as the first slide
\end{frame}

%\begin{frame}
%\frametitle{Overview} % Table of contents slide, comment this block out to remove it
%\tableofcontents % Throughout your presentation, if you choose to use \section{} and \subsection{} commands, these will automatically be printed on this slide as an overview of your presentation
%\end{frame}

%------------------------------------------------------------------------
% PRESENTATION SLIDES
%------------------------------------------------------------------------
\section{Where we landed}

\begin{frame}[t]{Yesterday: the average effect, estimated without bias}
\vfill
\begin{itemize} \vfill
\item \textbf{Estimand}: the ATE $\bar{\tau} = \bar{y}(1) - \bar{y}(0)$ --- a \mh{composite}, average claim \vfill
\item \textbf{Estimator}: Difference-in-Means $\hat{\tau}\left(\bm{Z}, \bm{Y}\right)$, treated mean $-$ control mean \vfill
\item Under CRA, \emph{whatever} the potential outcomes, $\E\left[\hat{\tau}\left(\bm{Z}, \bm{Y}\right)\right] = \bar{\tau}$ \vfill
\item ACORN study: the Difference-in-Means produced $\hat{\tau}\left(\bm{Z}, \bm{Y}\right) \approx 0.036$ \vfill
\item True variance (Neyman):
\begin{equation*}
\Var\left[\hat{\tau}\left(\bm{Z}, \bm{Y}\right)\right] = \frac{S^2_{\bm{y}(\bm{1})}}{n_1} + \frac{S^2_{\bm{y}(\bm{0})}}{n_0} - \frac{S^2_{\bm{\tau}}}{N}
\end{equation*} \vfill
\item $S^2_{\bm{\tau}}$: \mh{no data} to estimate it --- no unit reveals both POs \vfill
\end{itemize} \vfill
\end{frame}
%------------------------------------------------------------------------
\begin{frame}[t]{Today: from true variance to honest inference}
\vfill
\begin{itemize} \vfill
\item One experiment, \emph{one} estimate --- how far might $\hat{\tau}\left(\bm{Z}, \bm{Y}\right)$ sit from $\bar{\tau}$? \vfill
\item \textbf{Problem}: the true variance needs $S^2_{\bm{\tau}}$ --- and no data can estimate it \vfill
\item Today: \mh{ACORN}, plus \mh{village heads} to illustrate \citep[Ch.~2]{gerbergreen2012}: \vfill
\begin{itemize} \vfill
\item Conservative \emph{estimator of variance} $\to$ standard error estimator \vfill
\item Normal approximation $\to$ test statistic, $p$-values \vfill
\item Confidence intervals as the nulls we do not reject \vfill
\end{itemize} \vfill
\end{itemize} \vfill
\end{frame}
%------------------------------------------------------------------------
\section{Conservative variance estimation}

\begin{frame}[t]{The variance-estimation problem}
\vfill
\begin{itemize} \vfill
\item True variance has three pieces:
\begin{equation*}
\Var\left[\hat{\tau}\left(\bm{Z}, \bm{Y}\right)\right] = \frac{S^2_{\bm{y}(\bm{1})}}{n_1} + \frac{S^2_{\bm{y}(\bm{0})}}{n_0} - \frac{S^2_{\bm{\tau}}}{N}
\end{equation*} \vfill
\item \mh{None} directly observable --- all are features of fixed POs \vfill
\item But treated units reveal a \mh{sample} of $y_i(1)$; control units of $y_i(0)$ \vfill
\item[] $\Rightarrow$ \emph{estimate} $S^2_{\bm{y}(\bm{1})}$ (treated arm) and $S^2_{\bm{y}(\bm{0})}$ (control arm) \vfill
\item No unit reveals \emph{both} POs $\Rightarrow$ no sample of effects $\tau_i$ \vfill
\item[] $\Rightarrow$ no data to estimate $S^2_{\bm{\tau}}$ \vfill
\end{itemize} \vfill
\end{frame}
%------------------------------------------------------------------------
\begin{frame}[t]{Neyman's conservative estimator of the variance}
\vfill
\begin{itemize} \vfill
\item Split off the inestimable term: $\Var\left[\hat{\tau}\right] = \underbrace{\dfrac{S^2_{\bm{y}(\bm{1})}}{n_1} + \dfrac{S^2_{\bm{y}(\bm{0})}}{n_0}}_{U} - \dfrac{S^2_{\bm{\tau}}}{N}$ \vfill
\item $S^2_{\bm{\tau}}/N \geq 0$ $\Rightarrow$ $U$ is an \mh{upper bound}: $U \geq \Var\left[\hat{\tau}\right]$ \vfill
\item We can \emph{unbiasedly} estimate $S^2_{\bm{y}(\bm{1})}, S^2_{\bm{y}(\bm{0})}$, so
\begin{equation*}
\widehat{\Var}\left[\hat{\tau}\right] = \frac{\hat{S}^2_{\bm{y}(\bm{1})}}{n_1} + \frac{\hat{S}^2_{\bm{y}(\bm{0})}}{n_0}, \qquad \E\left[\widehat{\Var}\right] = U \geq \Var\left[\hat{\tau}\right]
\end{equation*} \vfill
\item Why \mh{conservative}? overstating $\Var$ $\Rightarrow$ wider intervals, harder to reject \vfill
\item[] $\Rightarrow$ a true null is rejected \emph{less} than $\alpha$ of the time --- not too \mh{liberal} \vfill
\end{itemize} \vfill
\end{frame}
%------------------------------------------------------------------------
\begin{frame}[t]{Exact when effects are constant, conservative otherwise}
\vfill
\begin{itemize} \vfill
\item The gap $\E\left[\widehat{\Var}\right] - \Var\left[\hat{\tau}\right]$ is exactly the dropped term $S^2_{\bm{\tau}}/N$ \vfill
\item \mh{Constant} (homogeneous) effect $\tau_i = \bar{\tau}$ for all $i$ $\Rightarrow$ $S^2_{\bm{\tau}} = 0$ $\Rightarrow$ \emph{exact} \vfill
\item \mh{Varying} effects $\Rightarrow$ $S^2_{\bm{\tau}} > 0$ $\Rightarrow$ overstates (conservative) \vfill
\item Either way it \emph{never understates} uncertainty --- an honest default \vfill
\end{itemize} \vfill
\end{frame}
%------------------------------------------------------------------------
\begin{frame}[t]{Village heads: a worked example (both POs known)}
\vfill
\begin{itemize}
\item $N = 7$, $n_1 = 2$; budget share (\%) \citep[Ch.~2]{gerbergreen2012} \vfill 
\item[] \vfill 
\begin{center}
\small
\begin{tabular}{l|rrr} \hline
Village &$\bm{y}(\bm{0})$& $\bm{y}(\bm{1})$& $\bm{\tau}$  \\ \hline
1& 10 & 15  & 5  \\
2& 15 & 15  & 0   \\
3& 20 & 30  & 10   \\
4& 20 & 15  & -5   \\
5& 10 & 20  & 10   \\
6& 15 & 15  & 0   \\
7& 15 & 30  & 15   \\ \hline
Average & 15 & 20 & 5  \\ \hline
\end{tabular}
\end{center}
\item[] \vfill
\item Both columns known $\Rightarrow$ we can compute the \mh{true} variance, $\bar{\tau} = 5$ \vfill
\end{itemize} \vfill
\end{frame}
%------------------------------------------------------------------------
\begin{frame}{Village heads: the conservative estimator overstates}
\begin{itemize}
\item Over the $21$ equally likely assignments: \vfill 
\begin{itemize} \vfill 
\item True $\Var \approx 21.2$ (solid) \vfill 
\item $\E[\widehat{\Var}] \approx 28.3$ (dashed)
\end{itemize}
\item[]
\begin{figure}[H]
\includegraphics[width=0.6\linewidth]{figures/villageheads_consv_var.pdf}
\end{figure}
\item \footnotesize Gap $= \E[\widehat{\Var}] - \Var = 28.3 - 21.2 = 7.1 = S^2_{\bm{\tau}}/N$ --- \emph{exactly} the dropped term \normalsize
\end{itemize}
\end{frame}
%------------------------------------------------------------------------
\begin{frame}[fragile]{ACORN in \texttt{R}: the conservative standard error}
\begin{lstlisting}[style=Rstyle, basicstyle=\ttfamily\tiny, frame=single]
# Conservative variance: estimate the two arm pieces; no data for -S2_tau/N, so drop it
conservative_var <- function(z, y) {
  var(y[z == 1]) / sum(z == 1) + var(y[z == 0]) / sum(z == 0)
}
var_hat <- conservative_var(z = z_obs, y = y_obs)   # ~ 0.000595
se_hat  <- sqrt(var_hat)                            # ~ 0.0244
\end{lstlisting}
\begin{itemize} \vfill
\item True se $= \sqrt{\Var\left[\hat{\tau}\right]}$ --- unknown; plug in $\widehat{\text{se}} = \sqrt{\widehat{\Var}} \approx 0.024$ \vfill
\item $\widehat{\Var}$ is \mh{conservative}: $\E[\widehat{\Var}] \geq \Var$ \vfill
\item But $\widehat{\text{se}} = \sqrt{\widehat{\Var}}$ is \emph{not} --- $\E[\widehat{\text{se}}] \leq \sqrt{\E[\widehat{\Var}]}$ (Jensen) \vfill
\item[] Still $\widehat{\text{se}} \overset{p}{\to} \sqrt{U} \geq$ true se: conservative in the \mh{limit} \vfill
\item Next: use $\widehat{\text{se}}$ to test hypotheses and build intervals for $\bar{\tau}$ \vfill
\end{itemize}
\end{frame}
%------------------------------------------------------------------------
\section{Normal approximation}

\begin{frame}[t]{Why we now need an approximation}
\vfill
\begin{itemize} \vfill
\item We want to \mh{test} hypotheses about the average effect $\bar{\tau}$ \vfill
\item Fisher's \mh{sharp} null fixed \emph{every} PO $\Rightarrow$ exact randomization distribution \vfill
\item A null on the \mh{mean} $\bar{\tau} = \tau_0$ fixes only the average \vfill
\item[] Each missing PO still unknown $\Rightarrow$ no exact distribution of $\hat{\tau}$ to compute \vfill
\item \mh{Finite population CLT} \citep{liding2017} \vfill 
\begin{itemize} \vfill 
\item A Normal approximation needing only mean and variance \vfill
\item Describes $\hat{\tau}$'s distribution \emph{without} the missing POs \vfill
\end{itemize} \vfill 
\end{itemize} \vfill
\end{frame}
%------------------------------------------------------------------------
\begin{frame}{What the Normal approximates}
\begin{itemize}
\item Standardize the Difference-in-Means by its \mh{true} mean and variance:
\begin{equation*}
\frac{\hat{\tau}\left(\bm{Z}, \bm{Y}\right) - \bar{\tau}}{\sqrt{\Var\left[\hat{\tau}\left(\bm{Z}, \bm{Y}\right)\right]}} \overset{\text{approx.}}{\sim} \mathcal{N}\left(0, 1\right)
\end{equation*}
\item Standardized $\Rightarrow$ mean $0$, sd $1$ at every $N$; only the \mh{shape} converges \vfill
\item[]
\begin{figure}[H]
\includegraphics[width=\linewidth]{figures/villageheads_normal_approx.pdf}
\end{figure}
\item[] \footnotesize Village heads, stacked: the standardized estimator approaches $\mathcal{N}(0,1)$ as $N$ grows \normalsize
\end{itemize}
\end{frame}
%------------------------------------------------------------------------
\section{Hypothesis testing}

\begin{frame}[t]{A test statistic for the average effect}
\vfill
\begin{itemize} \vfill
\item Under $H_0: \bar{\tau} = \tau_0$, the estimator is centered there: $\E[\hat{\tau}] = \bar{\tau} = \tau_0$ \vfill
\item So measure the gap $\hat{\tau} - \tau_0$ in \mh{estimated} standard-error units:
\begin{equation*}
T = \frac{\hat{\tau}\left(\bm{Z}, \bm{Y}\right) - \tau_0}{\widehat{\text{se}}}, \qquad \widehat{\text{se}} = \sqrt{\widehat{\Var}\left[\hat{\tau}\right]}
\end{equation*} \vfill
\item Under $H_0$, the \mh{reference} is the standard Normal (finite-population CLT) \vfill
\item $\lvert T \rvert$ large $=$ $\hat{\tau}$ many estimated se's from $\tau_0$ \vfill
\item[] $\widehat{\text{se}}$ only \emph{conservative in the limit} ($\overset{p}{\to} \sqrt{U} \geq$ true se) \vfill
\end{itemize} \vfill
\end{frame}
%------------------------------------------------------------------------
\begin{frame}{Rejection region: alternative of a smaller ATE}
\begin{itemize}
\item $T$ in the lower tail $\Rightarrow$ reject $H_0$ in favor of a smaller ATE \vfill
\item[]
\begin{figure}[H]
\includegraphics[width=0.74\linewidth]{figures/stand_norm_plot_lower_tail.pdf}
\end{figure}
\end{itemize}
\end{frame}
%------------------------------------------------------------------------
\begin{frame}{Rejection region: alternative of a larger ATE}
\begin{itemize}
\item $T$ in the upper tail $\Rightarrow$ reject $H_0$ in favor of a larger ATE \vfill
\item[]
\begin{figure}[H]
\includegraphics[width=0.74\linewidth]{figures/stand_norm_plot_upper_tail.pdf}
\end{figure}
\end{itemize}
\end{frame}
%------------------------------------------------------------------------
\begin{frame}{Rejection region: two-sided alternative}
\begin{itemize}
\item $T$ in \emph{either} tail $\Rightarrow$ reject $H_0$; split $\alpha$ across the two tails \vfill
\item[]
\begin{figure}[H]
\includegraphics[width=0.74\linewidth]{figures/stand_norm_plot_two_side.pdf}
\end{figure}
\end{itemize}
\end{frame}
%------------------------------------------------------------------------
\begin{frame}[t]{$p$-values: a tail area under the standard Normal}
\vfill
\begin{itemize} \vfill
\item \mh{$p$-value}: chance of a $T$ at least as extreme as observed, if $H_0$ true \vfill
\item A \mh{tail area} of the standard Normal, beyond the observed $T$ \vfill
\item $\Phi$ = standard Normal CDF; $1 - \Phi(T)$ = upper-tail area \vfill
\item[]
\begin{equation*}
p_{\text{upper}} = 1 - \Phi(T), \qquad p_{\text{lower}} = \Phi(T), \qquad p_{\text{two}} = 2\left(1 - \Phi\left(\lvert T \rvert\right)\right)
\end{equation*} \vfill
\item Each tail matches its \mh{alternative}: upper, lower, or two-sided \vfill
\item \mh{Decision}: reject $H_0$ when $p \leq \alpha$ \vfill
\end{itemize}
\end{frame}
%------------------------------------------------------------------------
\begin{frame}{Why the test is valid}
\begin{itemize}
\item \mh{Valid}: $\Pr\left(\text{reject } H_0 \mid H_0 \text{ true}\right) \leq \alpha$ \vfill
\item Write $T$ as two factors; their limits combine (Slutsky's Theorem):
\begin{equation*}
T = \underbrace{\frac{\hat{\tau} - \tau_0}{\sqrt{\Var\left[\hat{\tau}\right]}}}_{\to\, \mathcal{N}(0,1)\ \text{under}\ H_0} \times \underbrace{\frac{\sqrt{\Var\left[\hat{\tau}\right]}}{\widehat{\text{se}}}}_{\overset{p}{\to}\, c\, \leq\, 1} \ \overset{d}{\to}\ \mathcal{N}(0, c^2)
\end{equation*}
\item Slutsky: $T$ converges to $cZ$, with $Z \sim \mathcal{N}(0,1)$, $c$ a constant \vfill
\begin{itemize}
\item Property of the \mh{limiting} $cZ$, not a finite-$N$ rescaling \vfill
\item $\E[cZ] = 0$: scaling \mh{fixes the center} at $0$ \vfill
\item $\Var[cZ] = c^2$: scaling \mh{shrinks the spread} \vfill
\end{itemize}
\item \mh{Why $c \leq 1$?} conservative $\widehat{\text{se}} \overset{p}{\to} \sqrt{U} \geq \sqrt{\Var[\hat{\tau}]}$ \vfill
\item Dividing by \emph{over-large} $\widehat{\text{se}}$ shrinks every $T$ toward $0$, hence \mh{smaller variance}
\end{itemize}
\end{frame}

\begin{frame}{Why the test is valid (cont.)}
\begin{figure}[H]
\includegraphics[width=0.8\linewidth]{figures/consv_stand_norm_plot.pdf}
\end{figure}
\begin{itemize}
\item[] \footnotesize Rejection region holds \emph{less} mass under the narrower law $\Rightarrow$ reject a true null with prob.\ $\leq \alpha$ \normalsize
\end{itemize}
\end{frame}
%------------------------------------------------------------------------
\begin{frame}[fragile]{ACORN in \texttt{R}: testing no average effect}
\begin{lstlisting}[style=Rstyle, basicstyle=\ttfamily\tiny, frame=single]
# test statistic for H0: ATE = 0  (null mean 0, conservative se)
T_obs <- (tau_hat - 0) / se_hat                  # ~ 1.49

# Normal-approximation p-values (Phi = pnorm)
pnorm(q = T_obs, lower.tail = FALSE)             # upper p    ~ 0.068
2 * pnorm(q = abs(T_obs), lower.tail = FALSE)    # two-sided  ~ 0.137
\end{lstlisting}
\begin{itemize} \vfill
\item At $\alpha = 0.05$: \mh{do not reject} $\bar{\tau} = 0$ (two-sided $p \approx 0.137$) \vfill
\item Estimate is positive, but \emph{not} distinguishable from $0$ at this size \vfill
\end{itemize}
\end{frame}
%------------------------------------------------------------------------
\section{Confidence intervals by inverting tests}

\begin{frame}{Confidence interval by inverting the test}
\begin{itemize}
\item \mh{Confidence interval}: the null ATEs $\tau_0$ we do \emph{not} reject \vfill
\item Test each $\tau_0$; keep it when $\lvert T(\tau_0) \rvert \leq 1.96$ \vfill
\item[]
\begin{figure}[H]
\includegraphics[width=0.62\linewidth]{figures/stand_norm_plot_two_side_CI.pdf}
\end{figure}
\item[] \footnotesize Endpoints are the $\tau_0$ giving $T = \pm 1.96$: ACORN $95\%$ CI $= [-0.011,\ 0.084]$ \normalsize
\end{itemize}
\end{frame}
%------------------------------------------------------------------------
\begin{frame}[t]{From inversion to the familiar formula}
\vfill
\begin{itemize} \vfill
\item $z_{1-\alpha/2}$: the standard-Normal $1 - \alpha/2$ quantile (e.g.\ $1.96$ when $\alpha = 0.05$) \vfill
\item Retain $\tau_0$ with $\left\lvert \dfrac{\hat{\tau}\left(\bm{Z}, \bm{Y}\right) - \tau_0}{\widehat{\text{se}}} \right\rvert \leq z_{1 - \alpha/2}$ \vfill
\item Absolute value $\leq z_{1-\alpha/2}$ means \emph{both} tails at once:
\begin{equation*}
-\,z_{1 - \alpha/2} \;\leq\; \frac{\hat{\tau} - \tau_0}{\widehat{\text{se}}} \;\leq\; z_{1 - \alpha/2}
\end{equation*} \vfill
\item Now isolate $\tau_0$ in the middle --- three steps \vfill
\end{itemize} \vfill
\end{frame}

\begin{frame}[t]{Solving the inequality for $\tau_0$}
\vfill
\begin{itemize} \vfill
\item Multiply each part by $\widehat{\text{se}} > 0$ (no flip):
\begin{equation*}
-\,z_{1 - \alpha/2}\, \widehat{\text{se}} \;\leq\; \hat{\tau}\left(\bm{Z}, \bm{Y}\right) - \tau_0 \;\leq\; z_{1 - \alpha/2}\, \widehat{\text{se}}
\end{equation*} \vfill
\item Subtract $\hat{\tau}\left(\bm{Z}, \bm{Y}\right)$ from each part:
\begin{equation*}
-\,\hat{\tau}\left(\bm{Z}, \bm{Y}\right) - z_{1 - \alpha/2}\, \widehat{\text{se}} \;\leq\; -\tau_0 \;\leq\; -\hat{\tau}\left(\bm{Z}, \bm{Y}\right) + z_{1 - \alpha/2}\, \widehat{\text{se}}
\end{equation*} \vfill
\item Multiply by $-1$ (\emph{flips} both inequalities):
\begin{equation*}
\hat{\tau}\left(\bm{Z}, \bm{Y}\right) - z_{1 - \alpha/2}\, \widehat{\text{se}} \;\leq\; \tau_0 \;\leq\; \hat{\tau}\left(\bm{Z}, \bm{Y}\right) + z_{1 - \alpha/2}\, \widehat{\text{se}}
\end{equation*} \vfill
\item So $\hat{\tau} \pm z_{1 - \alpha/2}\, \widehat{\text{se}}$ is a \mh{consequence} of inversion, not a definition \vfill
\item ACORN $95\%$: $0.036 \pm 1.96(0.024) = [-0.011,\ 0.084]$, includes $0$ \vfill
\end{itemize} \vfill
\end{frame}
%------------------------------------------------------------------------
\section{Recap}

\begin{frame}[t]{Recap: from design to inference}
\vfill
\begin{itemize} \vfill
\item Yesterday: \mh{unbiased} $\bar{\tau}$, and its true variance (Neyman) \vfill
\item Today: a \mh{conservative} variance estimator $\Rightarrow$ honest standard error \vfill
\item \mh{Normal approximation} $\Rightarrow$ test statistic and $p$-values \vfill
\item \mh{Confidence interval} $=$ nulls we do not reject (inversion) \vfill
\item[] next: \mh{covariance adjustment} --- precision from baseline covariates \vfill
\end{itemize} \vfill
\end{frame}
%------------------------------------------------------------------------

\begin{frame}[allowframebreaks]
\frametitle{References}
\scriptsize
\bibliographystyle{chicago}
\bibliography{Master_Bibliography}
\end{frame}
%------------------------------------------------------------------------
\end{document}
